% TOP_PROBLEM: {"id": 1159, "title": "The Gauss Circle Problem", "category_id": 1, "category": {"id": 1, "name": "number_theory", "display_name": "Number Theory", "description": "Properties of integers, prime numbers, Diophantine equations.", "slug": "number-theory", "order_index": 1, "created_at": "2026-07-31T15:26:25.670Z"}, "status": "open"}
% FIRST_POSTED: 2026-09-25
% ENTRY_KIND: statement_only
% !TeX program = lualatex
% Definition and sources only; this file is not an LLM solution attempt.
\documentclass[11pt]{article}
\usepackage[margin=25mm]{geometry}
\usepackage{fontspec}
\setmainfont{Latin Modern Roman}
\usepackage{amsmath,amssymb,mathtools}
\usepackage{unicode-math}
\setmathfont{Latin Modern Math}
\usepackage{xurl,hyperref}
\hypersetup{colorlinks=true,urlcolor=blue}
\setcounter{secnumdepth}{0}
\setlength{\parindent}{0pt}
\setlength{\parskip}{5pt}
\setlength{\emergencystretch}{3em}
\begin{document}
\section*{\texorpdfstring{The Gauss Circle Problem}{The Gauss Circle Problem}}\label{1159}
\subsection{Definitions and mathematical statement}
For $r\ge0$, define
\[
 N(r)=\#\{(m,n)\in{\mathbb{Z}}^2:m^2+n^2\le r^2\},\qquad E(r)=N(r)-\pi r^2.
\]
The main selected estimate asks whether, for every ${\varepsilon}>0$, there are constants $C_{\varepsilon},r_{\varepsilon}$ with
$|E(r)|\le C_{\varepsilon} r^{1/2+{\varepsilon}}$ for all $r\ge r_{\varepsilon}$.
The broader request is the sharp order of this lattice-counting error. Radius, not squared radius, is the variable; changing variables changes the numerical exponent. The conjectured estimate does not assert the endpoint bound $O(r^{1/2})$ or a fixed sign of the discrepancy.
\par\smallskip\textit{Formulation and concept references:} \href{https://mathworld.wolfram.com/GausssCircleProblem.html}{[S1]}.

\subsection{Short English statement}
How closely does the number of lattice points in a large disk track its area, and is the error bounded by every power just above the square root of its radius?
\subsection{Sources}
\begin{itemize}
\item[S1]Formal statement, Bounds on a solution and conjecture.
\url{https://mathworld.wolfram.com/GausssCircleProblem.html}.
\textit{Formulation source.}
\end{itemize}
\end{document}
